% SOURCE: Chinese PDF viewer pp.133--134 (first half).
\section{Born Series}
\subsection*{1) First Born approximation}
For the integral equation and scattering amplitude,
\begin{equation}
\psi^+(\vect r)\xrightarrow{r\to\infty}
Ae^{\ii\vect k\cdot\vect r}
-\frac1{4\pi}\frac{2m}{\hbar^2}\frac{e^{\ii kr}}r
\int\dd^3r'\,e^{-\ii\vect k'\cdot\vect r'}
V(\vect r')\psi^+(\vect r'),
\end{equation}
\begin{equation}
f(\theta,\varphi)=-\frac{(2\pi)^3}{4\pi}\frac{2m}{\hbar^2}
\int\dd^3r'\,\frac{e^{-\ii\vect k'\cdot\vect r'}}{(2\pi)^{3/2}}
V(\vect r')\psi^+(\vect r'),
\end{equation}
if the potential $V(\vect r)$ is weak (incident energy $E\gg V$), one may solve iteratively. The zeroth-order approximation is
\begin{equation}
\psi^{(0)}(\vect r)=\frac{e^{\ii\vect k\cdot\vect r}}{(2\pi)^{3/2}}.
\end{equation}
The first Born approximation then gives
\begin{equation}\begin{aligned}
f(\theta,\varphi)
&=-\frac{(2\pi)^3}{4\pi}\frac{2m}{\hbar^2}
\int\dd^3r'\,\frac{e^{-\ii\vect k'\cdot\vect r'}}{(2\pi)^{3/2}}
V(\vect r')\frac{e^{\ii\vect k\cdot\vect r'}}{(2\pi)^{3/2}}\\
&=-\frac1{4\pi}\int\dd^3r'\,
 e^{\ii(\vect k-\vect k')\cdot\vect r'}U(\vect r').
\end{aligned}\end{equation}
For low-energy scattering (still satisfying $E\gg V$),
\begin{equation}\begin{aligned}
q&=|\vect k-\vect k'|=
\sqrt{k^2+k'^2-2kk'\cos\theta}\\
&=\sqrt{2k^2(1-\cos\theta)}
=\sqrt{4k^2\sin^2\frac\theta2}=2k\sin\frac\theta2.
\end{aligned}\end{equation}
If $q$ is small, the phase may be regarded as constant over the effective region of the interaction, and contributes no change to the cross section. Hence
\begin{equation}
\boxed{f(\theta,\varphi)=-\frac1{4\pi}\int\dd^3r'\,U(\vect r').}
\end{equation}

\textbf{Example.} Consider low-energy scattering by
\begin{equation}
V(r)=\begin{cases}V_0,&r\le a,\\0,&r>a.\end{cases}
\end{equation}
Then
\begin{equation}\begin{aligned}
f(\theta,\varphi)
&=-\frac1{4\pi}\int\dd^3r\,U(r)
=-\frac{m}{2\pi\hbar^2}\int\dd^3r\,V(r)\\
&=-\frac{mV_0}{2\pi\hbar^2}\int_{r\le a}\dd^3r
=-\frac{mV_0}{2\pi\hbar^2}\frac{4\pi}{3}a^3.
\end{aligned}\end{equation}
The differential scattering cross section is
\begin{equation}
\sigma(\theta,\varphi)=|f(\theta,\varphi)|^2
=\left(\frac{2mV_0a^3}{3\hbar^2}\right)^2.
\end{equation}
The total cross section is
\begin{equation}
\sigma_{\mathrm{tot}}=\int\sigma(\theta,\varphi)\dd\Omega
=4\pi\left(\frac{2mV_0a^3}{3\hbar^2}\right)^2.
\end{equation}

For general scattering by a central potential, $V(\vect r)=V(r)$,
\begin{equation}\begin{aligned}
f(\theta,\varphi)
&=-\frac1{4\pi}\int_0^\infty\dd r'\,r'^2U(r')
\int_0^\pi\dd\theta'\,\sin\theta' e^{-\ii qr'\cos\theta'}
\int_0^{2\pi}\dd\varphi'\\
&=-\frac1q\int_0^\infty\dd r'\,U(r')r'\sin(qr')
=-\frac{2m}{\hbar^2q}\int_0^\infty\dd r'\,V(r')r'\sin(qr').
\end{aligned}\end{equation}
Note: The Born approximation applies to scattering by a weak potential.

\subsection*{2) Born series}
The coordinate-space Lippmann--Schwinger equation is
\begin{equation}
\psi(\vect r)=\psi^{(0)}(\vect r)
+\int\dd^3r'\,\psi(\vect r')U(\vect r')G(\vect r,\vect r').
\end{equation}
\textbf{Zeroth-order approximation:}
\begin{equation}
\psi^{(0)}(\vect r)=Ae^{\ii\vect k\cdot\vect r}.
\end{equation}
\textbf{First-order approximation:}
\begin{equation}
\psi^{(1)}(\vect r)=\psi^{(0)}(\vect r)
+\int\dd^3r'\,\psi^{(0)}(\vect r')U(\vect r')G(\vect r,\vect r').
\end{equation}
\textbf{Second-order approximation:}
\begin{equation}\begin{aligned}
\psi^{(2)}(\vect r)
&=\psi^{(0)}(\vect r)
+\int\dd^3r'\,\psi^{(1)}(\vect r')U(\vect r')G(\vect r,\vect r')\\
&=\psi^{(0)}(\vect r)
+\int\dd^3r'\,\psi^{(0)}(\vect r')U(\vect r')G(\vect r,\vect r')\\
&\quad+\int\dd^3r'\dd^3r''\,\psi^{(0)}(\vect r')U(\vect r')
G(\vect r'',\vect r')U(\vect r'')G(\vect r,\vect r'').
\end{aligned}\end{equation}
The above series may be represented by diagrams:
\begin{center}
\begin{tikzpicture}[>=Latex,scale=.78,every node/.style={font=\scriptsize}]
 \begin{scope}[shift={(0,0)}]
 \draw[->] (-2.5,0)--(2.0,0) node[right] {$\psi_0$};
 \end{scope}
 \node at (3.38,0) {$+$};
 \begin{scope}[shift={(5.2,0)}]
 \draw[->] (-1.9,0)--(-.65,0) node[above] {$\psi_0$};
 \draw[thick] (0,0) circle (.7);\fill (0,0) circle (.07);\node[below] at (0,-.14) {$V$};
 \draw[->] (.65,0)--(1.9,0) node[above] {$\psi$};\node[above] at (.3,.89){$G$};
 \end{scope}
 \node at (8.0,0) {$+$};
 \begin{scope}[shift={(11.1,0)}]
 \draw[->] (-2.0,0)--(-1.5,0)node[above]{$\psi_0$};
 \draw[thick] (-1,0) circle (.53);\draw[thick] (.75,0) circle (.53);
 \fill (-1,0)circle(.07);\fill (.75,0)circle(.07);
 \draw[->](-.45,0)--(.18,0);\draw[->](1.3,0)--(2,0) node[above]{$\psi$};
 \node[below]at (-1,-.12){$V$};\node[below]at (.75,-.12){$V$};
 \node[above] at (-.1,.36){$G$};
 \end{scope}
\end{tikzpicture}
\end{center}
The Green function is the function that propagates interactions between two points.